\section{Discussion and open mathematical problems}
The positive result is a certified connection between an exact local optimizer and a complete feasible hybrid trajectory. Root isolation identifies an actual solution of the active equations. The gradient simplex proves strict local primary optimality in five flare coordinates. Stationary-time arguments and interval coverage establish feasibility, including the approach, occupied geometry and ideal contact law. The dependence structure then transfers the upper-support argument to a non-strict local maximum in twelve parameters.

The remaining global gap is concrete. The achieved lower endpoint in \eqref{eq:bracket} and the family exclusion differ by approximately $0.0296094506981383$. Closing that gap requires either a better feasible construction or a stronger exclusion using constraints discarded by the energy relaxation. The six active flare constraints suggest a lower-dimensional global search, but local sharpness does not exclude another branch. Nor does the terminal relaxation identify the controls needed to attain its scalar bound.

The secondary problem is also distinct. Continuous primary retention has been certified, and numerical effort minimizations have been delivered, but no inequality proves their effort optimality. A complete condition map would require validated continuation of active branches, treatment of changing active sets and geometric disjunctions, and either constructions or exclusions for the unresolved wind and urban cases. The three small wind neighborhoods establish only the explicitly parameterized feasible family.

The mathematical and physical domains should be kept separate. This reduced model omits calibrated aerodynamic uncertainty, six-degree-of-freedom attitude dynamics, load-rated contact, causal gust response and go-around viability. Its scene completes missing obstacle coverage through explicit assumptions. Those limitations do not negate the local theorem; they identify the object to which it applies. A historical or operational claim would require a different evidential and modelling argument.

